\section{Introduction}\label{sec:introduction}

Let $p$ be a complex polynomial of degree $n\ge2$, and let $Z(p)$
denote its zeros, counted with algebraic multiplicity. The Gauss--Lucas
theorem places the critical points $Z(p')$ in the convex hull
$\conv Z(p)$. Schmeisser's conjecture asserts that, when
$Z(p)\subseteq\overline\D=\{z\in\C:|z|\le1\}$, the closed unit disks
centered at the critical points cover this entire convex hull.
We prove the following degree range.

\begin{theorem}\label{thm:main}
Let $p$ be a complex polynomial of degree $n\ge6251$, all of whose
zeros lie in $\overline\D$. Then
\[
 \conv Z(p)\subseteq
 \bigcup_{\xi\in Z(p')}\{z\in\C:|z-\xi|\le1\}.
\]
\end{theorem}

The radius is sharp in every degree: for $p(z)=z^n-1$, the only
critical point is zero, while every zero has modulus one.

\subsection{The conjecture and its early development}

Sendov's conjecture concerns the same distance bound at each zero of
$p$. Sendov proposed it in 1958; Hayman's 1967 collection included it
as Problem~4.5 under Ilieff's name, through whom the problem had been
communicated \cite{Hayman1967,Schmeisser1971,Schmeisser2003}.
Rubinstein \cite{Rubinstein1968} proved the assertion at zeros on the
unit circle and for polynomials of degrees three and four.
Meir and Sharma \cite{MeirSharma1969} subsequently treated degree five.

A stronger localization already appeared in Schmeisser's early work.
For degrees at most four, he proved that for every zero $a$ there is
a critical point $\xi$ satisfying
\begin{equation}\label{in:early-localization}
             |\xi-a/2|\le1-|a|/2
\end{equation}
\cite{Schmeisser1969}. Gacs extended this localization to degree five
\cite{Gacs1971}. It implies the convex-hull covering directly.
Indeed, if $w=\sum_i t_i z_i$, with $t_i\ge0$ and $\sum_i t_i=1$,
then
\[
 \sum_i t_i\Re(w\overline{z_i})=|w|^2.
\]
Choose a zero $a$ for which $\Re(w\overline a)\ge|w|^2$. Then
\[
 |w-a/2|^2=|w|^2-\Re(w\overline a)+|a|^2/4
 \le|a|^2/4.
\]
Together with \eqref{in:early-localization}, this gives $|w-\xi|\le1$.

Schmeisser also established geometric and algebraic special cases.
His 1969 paper treats hulls whose vertices lie on the unit circle;
interior zeros are allowed \cite{Schmeisser1969}. His 1971 paper
develops coefficient estimates and proves Sendov's conjecture for
trinomials and for monic polynomials with real nonpositive
nonleading coefficients \cite{Schmeisser1971}. In Section~3 of
\cite{Schmeisser1977}, he formulated the convex-hull conjecture and
observed that the earlier proofs give this stronger conclusion in
degrees at most five, for hulls with all vertices on the unit circle,
and for trinomials. The same paper proves the assertion for triangular
hulls of arbitrary degree. Problem~4.30 of Hayman and Lingham
\cite{HaymanLingham} records the convex-hull question and its special
cases $p(0)=0$ and monic polynomials with real nonpositive
nonleading coefficients.

\subsection{Subsequent distance estimates}

For Sendov's conjecture, Brown and Xiang \cite{BrownXiang1999}
proved the result in degrees at most eight, and in arbitrary degree
when there are at most eight distinct zeros. The large-degree
approach began with D\'egot's theorem \cite{Degot}: for every fixed
$a\in(0,1)$ there is a degree bound depending on $a$ such that
a polynomial with $p(a)=0$ and all zeros in $\overline\D$ has
a critical point within distance one of $a$. Chalebgwa
\cite{Chalebgwa} made this dependence explicit, obtaining uniform
bounds when $a$ ranges over a compact subinterval of $(0,1)$.
Tao \cite{Tao2022} proved a single absolute, unspecified degree
threshold valid for all zeros. His argument treats separately
sequences of distinguished zeros whose moduli tend to zero,
to an interior point of $(0,1)$, or to one. An explicit absolute
bound $n\ge10^{200000}$ was subsequently given by Zhang
\cite{ZhangExplicit}. Mazur's all-degree argument and Tao's
exposition concern the distance from each actual zero to the
critical points \cite{Mazur,Tao2026}; the convex-hull assertion
appears as Conjecture~18 in \cite{Tao2026}.

A related line of work studies the centroid
\[
 G=\frac1n\sum_{z\in Z(p)}z
   =\frac1{n-1}\sum_{\xi\in Z(p')}\xi.
\]
Paw{\l}owski \cite{Pawlowski} proved, for zeros in $\overline\D$,
\[
 \min_{\xi\in Z(p')}|G-\xi|
 \le\frac{2n^{1/(n-1)}}{n^{2/(n-1)}+1}<1.
\]
Zhang improved this bound to $\sqrt{(n-2)/(n-1)}$
\cite{ZhangCentroid}. For $Z(p)\subseteq\overline\D$, his later paper
\cite{ZhangHull} gives
\[
 \min_{\xi\in Z(p')}|a-\xi|\le\sqrt{1-|a|^2}
       \qquad(a\in\conv Z(p')).
\]
The domain in this estimate is the critical-point hull.
For example, its radius follows from a convex representation
$a=\sum_j t_j\xi_j$:
\[
 \sum_j t_j|\xi_j-a|^2
 =\sum_j t_j|\xi_j|^2-|a|^2\le1-|a|^2.
\]
Schmeisser's question concerns the larger set $\conv Z(p)$.

Two recent estimates at an actual zero are useful here.
For zeros in $\overline\D$, Zhang \cite{ZhangQuadratic} proved the
quadratic Tang--Zhang inequality
\[
 \sum_{\xi\in Z(p')}\frac1{|a-\xi|^2}\ge n-1
       \qquad(a\in Z(p)),
\]
with an infinite summand at a common zero of $p$ and $p'$.
For arbitrary complex polynomials, Bouthat and Fallat
\cite{BouthatFallat}, and Zhang
\cite{ZhangVariance}, proved Borcea's quadratic variance bound
\[
 \min_{\xi\in Z(p')}|a-\xi|\le\sigma_2(p)
       \qquad(a\in Z(p)),\qquad
 \sigma_2(p)^2=\frac1n\sum_{z\in Z(p)}|z-G|^2.
\]
Section~\ref{sec:actual-roots} proves these estimates for
$n\ge6251$ and $n\ge2048$, respectively.

\subsection{The explicit degree bound}

An earlier argument established Schmeisser's conjecture for all
sufficiently large degrees \cite{TangQualitative}.
Theorem~\ref{thm:main} gives the explicit threshold $6251$.
Its proof consists of finite polynomial identities, analytic
inequalities uniform for $n\ge6251$, and a finite collection of
rational and interval verifications.

A point missed by all the critical-point unit disks can be moved
along a ray to a supporting edge of the zero hull. At a rightmost
zero $b$ we study the reciprocal critical points $(b-\xi)^{-1}$.
A matrix representation bounds their first two moments. A saddle
estimate and local primitive estimates sharpen these bounds and
produce nearby zeros of the original polynomial. Near the origin,
a finite Cayley--Hermite--Biehler identity supplies a source estimate;
the radial identity couples it to the same polynomial's moments.
In the other regions, equal-value fibers, product estimates and
positive sums of root counts yield a contradiction.
Section~\ref{sec:assembly} combines the parameter ranges.

The paper gives the analytic estimates, the complete remainders,
the parameter domains and the finite verification criteria.
The electronic appendix supplies the rational data and verification
programs. All zeros and critical points in sums and products are
counted with multiplicity. Terminating decimals denote exact rational
numbers. Each certified strict inequality holds on the indicated
closed parameter cell, including its boundary.
